\section{Анализ задачи}

\subsection{Проблема ручного маппинга с дополнительными параметрами}

В реальных проектах часто возникает ситуация, когда две модели данных имеют не полностью совпадающие наборы полей. Например, класс \texttt{A} содержит поля \texttt{x} и \texttt{y}, а класс \texttt{B} — поля \texttt{x} и \texttt{z}. При преобразовании \texttt{A → B} поле \texttt{z} не может быть получено из \texttt{A}, поэтому оно должно быть передано как дополнительный параметр функции маппера. Аналогично при обратном преобразовании \texttt{B → A} поле \texttt{y} является недостающим. Ручное написание таких функций требует перечисления всех совпадающих полей (их может быть десятки) и добавления параметров для отсутствующих. Это не только утомительно, но и чревато ошибками: легко пропустить какое-либо поле или перепутать тип. Кроме того, при изменении структуры моделей (добавлении, удалении или переименовании поля) все мапперы, использующие эти модели, нужно обновлять вручную, что делает поддержку кода дорогостоящей \cite{reff, refl}. Именно здесь метапрограммирование — автоматическая генерация кода — может существенно облегчить жизнь разработчика.

\subsection{Существующие решения и их недостатки}

Ниже рассмотрены основные библиотеки для автоматического маппинга в экосистеме Kotlin и Java, а также их недостатки применительно к Android-разработке \cite{refb, refc} и задачам генерации мапперов с дополнительными параметрами.

\textbf{ModelMapper (Java).} Эта библиотека использует рефлексию в рантайме для анализа структуры классов и автоматического копирования полей по именам. Разработчику достаточно вызвать \texttt{modelMapper.map(source, Target::class.java)}. Однако на Android рефлексия имеет серьёзные недостатки: она медленная (что критично для UI-потока), требует сложной настройки правил ProGuard/R8 для сохранения имён полей, и не работает с Kotlin-специфичными конструкциями, такими как data-классы с nullable-полями и inline-классы \cite{refa}. Кроме того, ModelMapper не поддерживает добавление произвольных параметров к функции маппинга — все поля целевого класса должны быть получены из исходного объекта. С точки зрения метапрограммирования, это пример \textit{рантайм-метапрограммирования} (рефлексии).

\textbf{MapStruct.} Популярная Java-библиотека, генерирующая реализации мапперов на этапе компиляции через APT. Она не использует рефлексию, поэтому работает быстро и безопасно с точки зрения обфускации. Это пример \textit{компиляционной генерации кода}. Однако для Kotlin-проектов MapStruct требует использования KAPT, который преобразует Kotlin-код в Java-стереотипы. Это приводит к значительному замедлению компиляции и потере информации о Kotlin-фичах: null-безопасность, перегрузки операторов, функции расширения обрабатываются некорректно. Также MapStruct не умеет генерировать мапперы, принимающие дополнительные параметры, которые не выводимы из исходного объекта.

\textbf{AutoMapper для Kotlin.} Существуют отдельные реализации AutoMapper, адаптированные для Kotlin, но большинство из них либо используют рефлексию через Kotlin-рефлексию, что ещё медленнее Java-рефлексии, либо работают через KAPT с теми же проблемами, что и MapStruct. Они также не предоставляют механизма для передачи дополнительных параметров в маппер.

\textbf{KSP-библиотеки (например, kotlin-mapper-processor).} Несколько проектов пытались использовать KSP для генерации мапперов, но они обычно решают только задачу «один-в-один» с одинаковыми именами полей, без поддержки вложенных объектов и коллекций, а также без возможности добавления внешних параметров. Кроме того, многие из них не поддерживают функции расширения и companion-объекты.

Таким образом, на рынке отсутствует решение, которое бы сочеталось с чистым Kotlin (без KAPT), поддерживало генерацию мапперов с дополнительными параметрами, обрабатывало вложенные структуры и коллекции, а также работало с функциями расширения и классами-объектами. Это обусловило необходимость разработки библиотеки \texttt{cast-castle}. Данный вывод перекликается с наблюдениями, сделанными в классических трудах по инженерии ПО: повторяющийся код ведёт к увеличению сложности и стоимости разработки \cite{refi, refl}.

\subsection{Требования к библиотеке cast-castle}

На основе проведённого анализа были сформулированы следующие требования к разрабатываемой библиотеке как инструменту метапрограммирования:

\begin{itemize}
    \item Аннотация \texttt{@CastCastleMapper} должна применяться к классам, функциям-членам и функциям расширения, а также к companion-объектам.
    \item Для аннотированного класса процессор должен генерировать реализации всех его абстрактных функций. Если функция имеет тело, но также аннотирована, процессор должен принудительно сгенерировать её перегрузку с добавленными параметрами.
    \item Для отдельной функции (standalone) процессор должен генерировать вспомогательную функцию с суффиксом \texttt{CastCastle}, принимающую дополнительные параметры — те поля целевого класса, которые отсутствуют в исходном типе.
    \item Для функции-расширения приёмник (\texttt{this}) считается исходным объектом; генерируемая функция должна быть также расширением или обычной функцией с параметром-приёмником.
    \item Имена сгенерированных функций образуются добавлением суффикса \texttt{CastCastle} к исходному имени (например, \texttt{map} → \texttt{mapCastCastle}).
    \item Процессор должен обрабатывать вложенные объекты: если поле исходного класса является другим классом, для которого также существует маппер, он должен быть вызван рекурсивно.
    \item Должна поддерживаться работа с коллекциями: \texttt{List}, \texttt{Set}, \texttt{MutableList} и их разновидности. Для каждого элемента коллекции должен применяться соответствующий маппер.
    \item Библиотека не должна использовать рефлексию в рантайме — вся генерация происходит на этапе компиляции (чистое метапрограммирование).
    \item Должна быть поддержка инкрементальной компиляции (incremental compilation) для ускорения пересборки.
\end{itemize}
